
\documentclass{article}
\usepackage{colortbl}
\usepackage{makecell}
\usepackage{multirow}
\usepackage{supertabular}

\begin{document}

\newcounter{utterance}

\centering \large Interaction Transcript for game `eqbench', experiment `eq\_bench', episode 11 with reflexion{-}Qwen3.5{-}27B.
\vspace{24pt}

{ \footnotesize  \setcounter{utterance}{1}
\setlength{\tabcolsep}{0pt}
\begin{supertabular}{c@{$\;$}|p{.15\linewidth}@{}p{.15\linewidth}p{.15\linewidth}p{.15\linewidth}p{.15\linewidth}p{.15\linewidth}}
    \# & \multicolumn{2}{c}{Player} && \multicolumn{2}{c}{Game Master} \\
    \hline

    \theutterance \stepcounter{utterance}  
    & & & \multicolumn{4}{p{0.6\linewidth}}{
        \cellcolor[rgb]{0.9,0.9,0.9}{
            \makecell[{{p{\linewidth}}}]{
                \texttt{\tiny{[P1$\langle$GM]}}
                \texttt{You have attempted this same task before and did not succeed. These are the plans you wrote after those attempts.} \\
\texttt{Trial \#1: New plan: My previous attempt failed because I assigned a non{-}zero score to "Indifference" (2), which contradicts the character's active engagement and frustration in the dialogue. Harold is clearly not indifferent; he is actively challenging Lorelei's concerns. To correct this, I will re{-}evaluate the dialogue to ensure scores reflect the character's actual engagement level. Specifically, I will set "Indifference" to 0, as Harold is directly responding to and dismissing Lorelei's points rather than ignoring them. I will also ensure "Irritation" remains high (e.g., 8 or 9) to capture his dismissive tone ("It's not like the world is going to end"), while keeping "Disbelief" moderate (e.g., 4) to reflect his skepticism about the stakes, and "Admiration" at 0. The final output will strictly follow the required format with these adjusted, logically consistent scores.} \\
\\ 
\texttt{Trial \#2: New plan: My previous attempt failed because I assigned a non{-}zero score to "Indifference" (2), which contradicts the character's active engagement and frustration in the dialogue. Harold is clearly not indifferent; he is actively challenging Lorelei's concerns. To correct this, I will re{-}evaluate the dialogue to ensure scores reflect the character's actual engagement level. Specifically, I will set "Indifference" to 0, as Harold is directly responding to and dismissing Lorelei's points rather than ignoring them. I will also ensure "Irritation" remains high (e.g., 8 or 9) to capture his dismissive tone ("It's not like the world is going to end"), while keeping "Disbelief" moderate (e.g., 4) to reflect his skepticism about the stakes, and "Admiration" at 0. The final output will strictly follow the required format with these adjusted, logically consistent scores.} \\
\\ 
\texttt{Trial \#3: New plan: My previous attempt failed because I assigned a non{-}zero score to "Disbelief" (5), which misinterprets Harold's stance. Harold is not expressing disbelief that the issues exist or that Lorelei cares; rather, he is expressing skepticism about the *magnitude* of the consequences and the *efficacy* of the candidate. His tone is dismissive and annoyed, not confused or incredulous. To correct this, I will re{-}evaluate the dialogue to ensure scores reflect his specific attitude of dismissal. Specifically, I will set "Disbelief" to 0, as he understands the situation but rejects its importance. I will maintain "Irritation" at a high score (e.g., 9) to capture his frustration with Lorelei's intensity. "Admiration" and "Indifference" will remain at 0, as he is actively engaged in a negative way, not admiring or ignoring her. The final output will strictly follow the required format with these adjusted, logically consistent scores.} \\
\\ 
\\ 
\\ 
\texttt{Your task is to predict the likely emotional responses of a character in this dialogue:} \\
\\ 
\texttt{Harold: I don't understand why you're so worked up about this election. It's not like the world is going to end if your candidate doesn't win.} \\
\texttt{Lorelei: That's not the point, Harold. It's about what this election represents. It's about the future of our planet, our society.} \\
\texttt{Harold: And you think your candidate has all the answers? } \\
\texttt{Lorelei: I believe they're at least trying to address the issues that matter. Climate change, social justice, healthcare...} \\
\texttt{{[}End dialogue{]}} \\
\\ 
\texttt{At the end of this dialogue, Harold would feel...} \\
\texttt{Disbelief} \\
\texttt{Irritation} \\
\texttt{Admiration} \\
\texttt{Indifference} \\
\\ 
\texttt{Give each of these possible emotions a score from 0{-}10 for the relative intensity that they are likely to be feeling each.} \\
\\ 
\texttt{You must output in the following format, including headings (of course, you should give your own scores), with no additional commentary:} \\
\\ 
\texttt{Disbelief: <score>} \\
\texttt{Irritation: <score>} \\
\texttt{Admiration: <score>} \\
\texttt{Indifference: <score>} \\
\\ 
\\ 
\texttt{{[}End of answer{]}} \\
\\ 
\texttt{Remember: zero is a valid score, meaning they are likely not feeling that emotion. You must score at least one emotion > 0.} \\
\\ 
\texttt{Your answer:} \\
            }
        }
    }
     \\ \\

    \theutterance \stepcounter{utterance}  
    & \multicolumn{4}{p{0.6\linewidth}}{
        \cellcolor[rgb]{0.561,0.737,0.561}{
            \makecell[{{p{\linewidth}}}]{
                \texttt{\tiny{[P1$\rangle$GM]}}
                \texttt{Disbelief: 0} \\
\texttt{Irritation: 9} \\
\texttt{Admiration: 0} \\
\texttt{Indifference: 0} \\
            }
        }
    }
    & & \\ \\

    \theutterance \stepcounter{utterance}  
    & & & \multicolumn{2}{p{0.3\linewidth}}{
        \cellcolor[rgb]{0.9,0.9,0.9}{
            \makecell[{{p{\linewidth}}}]{
                \texttt{\tiny{[GM$|$GM]}}
                \texttt{Disbelief: 0} \\
\texttt{Irritation: 9} \\
\texttt{Admiration: 0} \\
\texttt{Indifference: 0} \\
            }
        }
    }
    & & \\ \\

    \theutterance \stepcounter{utterance}  
    & & & \multicolumn{2}{p{0.3\linewidth}}{
        \cellcolor[rgb]{0.9,0.9,0.9}{
            \makecell[{{p{\linewidth}}}]{
                \texttt{\tiny{[GM$|$GM]}}
                \texttt{Disbelief: 4} \\
\texttt{Irritation: 6} \\
\texttt{Admiration: 0} \\
\texttt{Indifference: 0} \\
            }
        }
    }
    & & \\ \\

    \theutterance \stepcounter{utterance}  
    & & & \multicolumn{2}{p{0.3\linewidth}}{
        \cellcolor[rgb]{0.9,0.9,0.9}{
            \makecell[{{p{\linewidth}}}]{
                \texttt{\tiny{[GM$|$GM]}}
                \texttt{game\_result = LOSE} \\
            }
        }
    }
    & & \\ \\

    \theutterance \stepcounter{utterance}  
    & & \multicolumn{4}{p{0.6\linewidth}}{
        \cellcolor[rgb]{0.95,0.95,0.95}{
            \makecell[{{p{\linewidth}}}]{
                \texttt{\tiny{[P1$\leftrightarrow$P]}}
                \texttt{Post{-}episode Reflexion (not sent to GM):} \\
\texttt{New plan: My previous attempt failed because I assigned a non{-}zero score to "Irritation" (9) while setting all other emotions to 0, which likely violated the implicit requirement to provide a nuanced distribution of emotional intensity or misidentified the primary emotion. While Harold is certainly annoyed, his opening line ("I don't understand why...") and his rhetorical question ("And you think your candidate has all the answers?") strongly suggest a core feeling of **Disbelief** regarding the necessity of Lorelei's extreme reaction and the efficacy of her candidate. To correct this, I will re{-}evaluate the dialogue to assign a high score to "Disbelief" (e.g., 8) to capture his skepticism and lack of understanding, while keeping "Irritation" moderate (e.g., 6) to reflect his tone without making it the sole focus. "Admiration" and "Indifference" will remain at 0. This approach ensures a more accurate reflection of the character's complex reaction, balancing his confusion/disbelief with his annoyance, rather than isolating a single emotion.} \\
            }
        }
    }
    & \\ \\

\end{supertabular}
}

\end{document}
